\section{Introduction}
Building material-flow assessment requires several linked transformations. Observed building attributes are translated into material quantities; inventories are combined with temporal assumptions to estimate future flows; released materials are assessed for circular pathways; and these pathways are translated into environmental consequences. These transformations are supported by different modelling traditions. Probabilistic material-intensity characterization, dynamic material-flow analysis and circularity--life-cycle assessment provide complementary capabilities, but they do not necessarily establish the same evidential or probabilistic connection between their outputs and inputs \citep{U03,U12,U19}.

The relevant distinction is between the presence of uncertainty analysis and the continuity of uncertainty information. A downstream model may accept an upstream central estimate while disregarding its uncertainty. Alternatively, it may introduce a new uncertainty distribution without specifying whether it represents inherited error, additional process uncertainty or a scenario assumption. In either case, uncertainty at the final output cannot automatically be interpreted as a complete account of upstream uncertainty. The preliminary evidence audit therefore separates transfer status, representation change and changes in uncertainty magnitude rather than treating them as interchangeable outcomes.

The analytical chain used in this review comprises five stages: S1, observation and inference of building attributes; S2, material characterization and stock estimation; S3, temporal stock dynamics and material release; S4, recoverability and circular routing; and S5, environmental assessment and decision consequences. These are organizing categories for the review, not universal software layers. Studies can span several stages, bypass a stage or address only one transformation. Their actual variable relationships must therefore be examined alongside their initial stage assignments.

Digital twins are investigated here as a possible architectural response to interface limitations. Their contribution cannot be inferred from terminology, three-dimensional representation or the coexistence of several models. The review instead examines explicit capabilities for entity linkage, evidence integration, state updating, model coupling and interpretation of dependent outputs. Whether these capabilities preserve uncertainty or improve decisions is an empirical question for extraction, not an assumed benefit of digitization.

The intended contribution is a synthesized architecture whose components and connections can be traced to evidence and documented gaps. The research questions are:
\begin{enumerate}
\item What uncertainty information is transferred, transformed or left unassessed at the interfaces between S1--S5?
\item Which domain-specific digital twin capabilities implement these interfaces, and what limitations are reported or remain unresolved?
\item Which interdisciplinary mechanisms can address these limitations, under what transfer conditions, and with what level of supporting evidence?
\item How can these findings be synthesized into an architecture that accommodates spatial, temporal, physical and decision contexts without claiming universal applicability?
\end{enumerate}
